% ==================== 1.1. INTRODUCCIÓN ====================

\section{Introducción}\label{sec:intro_estado_arte}

Las enfermedades neurodegenerativas representan uno de los principales desafíos clínicos contemporáneos debido a su profundo impacto socioeconómico y a su progresión clínica irreversible. Patologías como la enfermedad de Parkinson, la enfermedad de Alzheimer o la Esclerosis Lateral Amiotrófica (ELA) comprometen severamente a una vasta proporción de la población mundial, induciendo un deterioro gradual de las capacidades motoras, cognitivas y lingüísticas \cite{abrilcarreres2024}. La detección precoz se ha consolidado como un factor clínico determinante, ya que la intervención temprana permite implementar estrategias terapéuticas más eficaces, ralentizar el declive funcional y optimizar la calidad de vida de los pacientes. Sin embargo, los protocolos diagnósticos tradicionales, fundamentados en evaluaciones neuropsicológicas o pruebas de neuroimagen, presentan limitaciones sustanciales: resultan invasivos, implican altos costes económicos y, frecuentemente, carecen de la sensibilidad necesaria para identificar alteraciones estructurales o funcionales en estadios iniciales \cite{avance2023}.

Como respuesta a estas limitaciones, el análisis computacional de la voz ha emergido como una alternativa diagnóstica altamente prometedora. La fonación constituye un proceso fisiológico complejo que exige la sincronización precisa de los subsistemas respiratorio, fonatorio, articulatorio y cognitivo. Consecuentemente, el daño tisular o la alteración neuromotora en el sistema nervioso central se proyecta directamente sobre las emisiones vocales \cite{romero2018}. Estas variaciones patológicas, a menudo imperceptibles para el sistema auditivo humano, pueden ser parametrizadas mediante técnicas de procesamiento digital de señales y empleadas como biomarcadores subclínicos. Su naturaleza intrínsecamente no invasiva, junto con su alta accesibilidad y bajo coste, sitúan a la señal de voz como un instrumento idóneo para el cribado poblacional y la monitorización continua del deterioro neurocognitivo.

De forma simultánea, la irrupción de la Inteligencia Artificial (IA) ha redefinido el análisis de bioseñales médicas. La aplicación de algoritmos de aprendizaje automático al estudio del habla ha posibilitado la creación de modelos predictivos con métricas de rendimiento que rivalizan o incluso superan a la instrumentación clínica estándar \cite{moro-velazquezUsingXVectorsAutomatically2020}. Históricamente, la literatura ha reportado la eficacia de métodos clásicos supervisados como Máquinas de Vectores de Soporte (SVM) y Bosques Aleatorios (\textit{Random Forests}), evolucionando posteriormente hacia arquitecturas de aprendizaje profundo (\textit{Deep Learning}) como Redes Neuronales Convolucionales (CNN) y Recurrentes (LSTM). En la actualidad, el paradigma experimenta una transición directa hacia el uso de \textit{foundation models} y técnicas de aprendizaje auto-supervisado (\textit{Self-Supervised Learning}, como Wav2Vec 2.0 o data2vec), arquitecturas preentrenadas sobre corpus masivos que extraen representaciones latentes del habla con una capacidad de generalización muy superior \cite{baevski2020, agbavorArtificialIntelligenceEnabledEndToEnd2023}.

Bajo esta premisa, el presente Estado del Arte persigue analizar y sintetizar la literatura científica más relevante en torno a la aplicación de la Inteligencia Artificial al estudio de la voz para la detección temprana de enfermedades neurodegenerativas. A lo largo del capítulo se examinarán de forma sistemática los enfoques metodológicos predominantes, las características acústicas de mayor valor diagnóstico y las métricas de desempeño reportadas en investigaciones contemporáneas. Asimismo, se identificarán las limitaciones técnicas actuales y se trazarán las líneas de investigación futuras, proporcionando el sustento teórico y metodológico necesario para el diseño, entrenamiento y evaluación de los modelos experimentales propuestos en este Trabajo de Fin de Grado.